% solutions/thm-a2-target.tex
% Standalone proof dossier for the A2 strong-Laplace transfer theorem.
\documentclass[../main.tex]{subfiles}
\begin{document}

% === SOLUTION HEADER (proof provenance) =====================================
%   ledger-node : thm:a2-target
%   refines     : thm:a2-target (ledger refines: conj:a2)
%   bounded_by  : obs:tv-insufficient; obs:gaussian-tail-rigidity
%   checked_by  : agent
%   author      : /root/a3_prior_prover
%   reviewer    : /root/a3_prior_review
%   review      : research/reviews/2026-08-27-a2-strong-laplace-post-promotion-proof-review.md
%   date        : 2026-08-27
% ============================================================================

\section*{Solution: strong-Laplace transfer of the three Gaussian constants}
\label{sec:sol-thm-a2-target}

\paragraph{Refined statement and probabilistic convention.}
All random objects below are functions of the data, whose law at sample size $n$ is denoted by
$\Prob_{\theta_0}$.  Thus $Y_n=o_P(1)$ means
$\Prob_{\theta_0}(\abs{Y_n}>\varepsilon)\to0$ for every $\varepsilon>0$; matrix convergence in
probability is in operator norm.  Values on events whose probability tends to zero may be defined
arbitrarily without changing any conclusion.

\begin{theorem}[Strong-Laplace form of Theorem~\ref{thm:a2-target}]
\label{thm:sol-a2-target}
Fix $d$, let $\theta_0$ be an interior data-generating truth, assume the prior density is positive
and smooth near $\theta_0$, and assume the model is regular.  For the resulting posterior
$\pi_n$, let $\hat\theta_n$ be its mode and $H_n$ its symmetric mode Hessian.  Suppose
\[
 \frac{H_n}{n}\ \xrightarrow{\ \Prob_{\theta_0}\ }\ J:=I(\theta_0)\succ0.
\]
On an event whose $\Prob_{\theta_0}$-probability tends to one, $H_n\succ0$ and the law
$\bar\pi_n$ of
\[
 z=H_n^{1/2}(\theta-\hat\theta_n),\qquad \theta\sim\pi_n,
\]
has density relative to $\gamma=N(0,\Id_d)$ given by
\[
 q_n(z):=\frac{\dd\bar\pi_n}{\dd\gamma}(z)
 =\frac{e^{-r_n(z)}}{\mathcal Z_n},
 \qquad
 \mathcal Z_n=\int e^{-r_n}\dd\gamma.
\]
Assume the global essential oscillation
\[
 \delta_n:=\operatorname*{ess\,sup}_{\gamma}r_n
          -\operatorname*{ess\,inf}_{\gamma}r_n
 =o_P(1).
\]
Then, with $\lambda_n=\lambda_{\max}(H_n^{-1})$,
\begin{equation}
\label{eq:sol-a2-relative-limits}
 \frac{\CP(\pi_n)}{\lambda_n},\qquad
 \frac{\CLS(\pi_n)}{\lambda_n},\qquad
 \frac{\CTCI(\pi_n)}{\lambda_n}
 \ \xrightarrow{\ \Prob_{\theta_0}\ }\ 1.
\end{equation}
Equivalently, each of the three constants is
$(1+o_P(1))\lambda_{\max}(H_n^{-1})$, and
\begin{equation}
\label{eq:sol-a2-fisher-limits}
 n\CP(\pi_n),\qquad n\CLS(\pi_n),\qquad n\CTCI(\pi_n)
 \ \xrightarrow{\ \Prob_{\theta_0}\ }\
 \lambda_{\max}\!\big(I(\theta_0)^{-1}\big).
\end{equation}
\end{theorem}

\begin{proof}
We first prove a deterministic comparison valid for each data realization on which $H_n\succ0$
and $\delta_n<\infty$.  Write
\[
 a_n=\operatorname*{ess\,inf}_{\gamma}r_n,\qquad
 b_n=\operatorname*{ess\,sup}_{\gamma}r_n,
 \qquad \delta_n=b_n-a_n.
\]
Since $\gamma$ is a probability measure,
\[
 e^{-b_n}\le\mathcal Z_n\le e^{-a_n}.
\]
Consequently the \emph{normalized} density ratio satisfies
\begin{equation}
\label{eq:sol-a2-normalized-ratio}
 e^{-\delta_n}\le q_n(z)\le e^{\delta_n}
 \quad\text{for $\gamma$-almost every }z.
\end{equation}
In particular, with $\eta_n=e^{\delta_n}-1$,
\begin{equation}
\label{eq:sol-a2-ratio-to-one}
 \norm{q_n-1}_{L^\infty(\gamma)}\le\eta_n.
\end{equation}

Let
\[
 T_n(z)=\hat\theta_n+H_n^{-1/2}z,\qquad
 G_n=(T_n)_\#\gamma=N(\hat\theta_n,H_n^{-1}).
\]
Then $\pi_n=(T_n)_\#\bar\pi_n$ and
\[
 \frac{\dd\pi_n}{\dd G_n}(\theta)
 =\frac{\exp\{-r_n(T_n^{-1}\theta)\}}{\mathcal Z_n}.
\]
The perturbing potential $r_n\circ T_n^{-1}$ has oscillation exactly $\delta_n$.
Under the repository normalization, the Gaussian calibration gives
\[
 \CP(G_n)=\CLS(G_n)=\lambda_n.
\]
The Holley--Stroock theorem therefore gives the sharp bounded-perturbation factor
\begin{equation}
\label{eq:sol-a2-hs-upper}
 \CP(\pi_n)\le e^{\delta_n}\lambda_n,\qquad
 \CLS(\pi_n)\le e^{\delta_n}\lambda_n.
\end{equation}
It is the oscillation $\delta_n$ of the unnormalized potential that appears here, not twice that
quantity from separately inserting the two loose bounds in
\eqref{eq:sol-a2-normalized-ratio}.  Otto--Villani, in the convention
$W_2^2\le2\CTCI\operatorname{KL}$, adds
\begin{equation}
\label{eq:sol-a2-tci-upper}
 \CTCI(\pi_n)\le\CLS(\pi_n)\le e^{\delta_n}\lambda_n.
\end{equation}

We next obtain a matching covariance lower bound.  Let $Z\sim\gamma$, and condition throughout
on the data realization.  Put
\[
 m_n=\int zq_n(z)\dd\gamma(z),\qquad
 M_n=\int zz^\top q_n(z)\dd\gamma(z),\qquad
 \Sigma_n=\Cov_{\bar\pi_n}(Z)=M_n-m_nm_n^\top.
\]
Because $\int z\dd\gamma=0$, \eqref{eq:sol-a2-ratio-to-one} and
$\E_\gamma\norm Z\le\sqrt d$ give
\begin{equation}
\label{eq:sol-a2-mean-bound}
 \norm{m_n}
 \le\eta_n\E_\gamma\norm Z
 \le\eta_n\sqrt d.
\end{equation}
For every unit vector $v$,
\[
 \abs{v^\top(M_n-\Id_d)v}
 =\abs{\int(q_n-1)(v^\top z)^2\dd\gamma(z)}
 \le\eta_n\int(v^\top z)^2\dd\gamma(z)=\eta_n.
\]
Hence
\begin{equation}
\label{eq:sol-a2-covariance-bound}
 \norm{\Sigma_n-\Id_d}_{\op}
 \le \eta_n+d\eta_n^2=:\varepsilon_n.
\end{equation}
Fixed dimension and $\delta_n=o_P(1)$ imply
$\eta_n=o_P(1)$ and $\varepsilon_n=o_P(1)$.  On the event $\varepsilon_n<1$,
\[
 (1-\varepsilon_n)\Id_d\preceq\Sigma_n
 \preceq(1+\varepsilon_n)\Id_d.
\]
Since covariance is translation invariant and transforms by congruence,
\[
 \Cov_{\pi_n}(\theta)
 =H_n^{-1/2}\Sigma_nH_n^{-1/2},
\]
and therefore
\begin{equation}
\label{eq:sol-a2-unwhitened-covariance}
 (1-\varepsilon_n)H_n^{-1}
 \preceq\Cov_{\pi_n}(\theta)
 \preceq(1+\varepsilon_n)H_n^{-1}.
\end{equation}
This relative Loewner comparison retains the anisotropy of $H_n$ and in particular gives
\begin{equation}
\label{eq:sol-a2-covariance-top}
 \lambda_{\max}(\Cov_{\pi_n})
 \ge(1-\varepsilon_n)\lambda_n.
\end{equation}

For the Poincar\'e constant, a linear test
$f_v(\theta)=v^\top\theta$ has energy $\norm{v}^2$ and variance
$v^\top\Cov_{\pi_n}v$.  The Gaussian density comparison ensures that these linear functions lie
in the form domain (or one may use smooth truncations).  Taking the supremum over $v\ne0$ gives
\[
 \CP(\pi_n)\ge\lambda_{\max}(\Cov_{\pi_n}).
\]
Linearization of LSI around constants then gives
\[
 \CLS(\pi_n)\ge\CP(\pi_n)\ge\lambda_{\max}(\Cov_{\pi_n}).
\]

For completeness, the transportation lower bound is separate; Poincar\'e and $T_2$ are not
ordered in general.  Suppose a probability measure $\mu$ with finite linear exponential moments
satisfies $T_2(C)$ in the repository convention.  For a linear function
$f_v(x)=v^\top x$, the entropy
variational formula and $W_1\le W_2$ imply, for every $t\in\R$,
\begin{align*}
 \log\E_\mu e^{t(f_v-\E_\mu f_v)}
 &=\sup_{\nu\ll\mu}
   \left\{t(\E_\nu f_v-\E_\mu f_v)-\KL(\nu\|\mu)\right\}\\
 &\le\sup_{k\ge0}
   \left\{\abs{t}\norm{v}\sqrt{2Ck}-k\right\}
 =\frac{Ct^2\norm{v}^2}{2}.
\end{align*}
Indeed, $T_2(C)$ gives
$W_1(\nu,\mu)\le W_2(\nu,\mu)\le\sqrt{2C\KL(\nu\|\mu)}$.
Expanding the centred log moment-generating function at $t=0$ yields
\[
 v^\top\Cov_\mu v\le C\norm{v}^2.
\]
Thus $T_2(C)$ implies $\Cov_\mu\preceq C\Id$, and taking optimal constants gives
\begin{equation}
\label{eq:sol-a2-tci-covariance-lower}
 \CTCI(\pi_n)\ge\lambda_{\max}(\Cov_{\pi_n}).
\end{equation}

Combining \eqref{eq:sol-a2-hs-upper}, \eqref{eq:sol-a2-tci-upper},
\eqref{eq:sol-a2-covariance-top}, and
\eqref{eq:sol-a2-tci-covariance-lower}, each
$C_n\in\{\CP(\pi_n),\CLS(\pi_n),\CTCI(\pi_n)\}$ obeys
\begin{equation}
\label{eq:sol-a2-final-sandwich}
 1-\varepsilon_n
 \le\frac{C_n}{\lambda_n}
 \le e^{\delta_n}
\end{equation}
whenever $\varepsilon_n<1$.  Both endpoints converge to one in
$\Prob_{\theta_0}$-probability, proving \eqref{eq:sol-a2-relative-limits}.

It remains to track the random matrix normalization rather than replace it informally by its
limit.  Set $B_n=H_n/n$.  Since $B_n\to J\succ0$ in operator norm in probability, the event
$B_n\succ0$ has probability tending to one, and continuity of inversion on the positive-definite
cone gives the matrix limit
\begin{equation}
\label{eq:sol-a2-matrix-limit}
 nH_n^{-1}=B_n^{-1}
 \ \xrightarrow{\ \Prob_{\theta_0}\ }\ J^{-1}
\end{equation}
in operator norm.  Continuity of the largest eigenvalue then yields
\[
 n\lambda_n=\lambda_{\max}(B_n^{-1})
 \ \xrightarrow{\ \Prob_{\theta_0}\ }\
 \lambda_{\max}(J^{-1}).
\]
Multiplying this scalar convergence by
\eqref{eq:sol-a2-relative-limits} and applying Slutsky's theorem proves
\eqref{eq:sol-a2-fisher-limits}.
\end{proof}

\paragraph{Hypotheses actually used.}
The transfer argument uses fixed $d$, the positive-definite matrix convergence
$H_n/n\to I(\theta_0)$, the global normalized density representation, and
$\operatorname{osc}(r_n)=o_P(1)$.  The location $\hat\theta_n$ matters only as a translation.
Once the matrix convergence and density representation are assumed, the interiority of
$\theta_0$, local prior smoothness, and model regularity are not invoked separately; in the
manuscript they are the statistical context intended to produce those two analytic inputs.
Growing dimension would require a quantitative replacement for
$\varepsilon_n=\eta_n+d\eta_n^2=o_P(1)$.

\paragraph{Obstructions respected.}
The formal fence \texttt{obs:tv-insufficient} is respected because
\eqref{eq:sol-a2-ratio-to-one} is global $L^\infty(\gamma)$ convergence of the whitened density
ratio, not total-variation convergence alone.  It controls all second moments and rules out the
remote-contamination mechanism used by that obstruction.  The second formal fence
\texttt{obs:gaussian-tail-rigidity} is also respected: the theorem asserts a Fisher-scale LSI
and $T_2$ only under global vanishing-oscillation comparison with the mode Gaussian; it makes no
such assertion for ordinary fixed-Gaussian-prior logistic posteriors, which exhibit the certified
prior-tail obstruction.

\paragraph{Certification boundary.}
There are no unclosed analytic steps and no numerical input.  The theorem is unconditional under
its explicit hypotheses; it has no unresolved ledger dependency.  The strength and
model-specific verifiability of the global oscillation hypothesis remain separate open questions.
This dossier is independently agent-certified; its mathematical scope ends at the stated
strong-Laplace implication and does not verify the global oscillation hypothesis in a model.

\end{document}
